\documentclass[12pt]{article}
\usepackage[letterpaper,margin=1in]{geometry}
\usepackage{amsmath,amssymb}
\usepackage{float}
\setlength{\abovecaptionskip}{0pt}\setlength{\belowcaptionskip}{6pt}
\usepackage[hidelinks]{hyperref}
\title{Magnus quadrature formulas of order 10 on a centered interval}
\author{Kenneth C. Johnson\\ \normalsize ken@kjinnovation.com}
\date{Version 02-Nov-2026}
\begin{document}
\maketitle

\section*{Abstract}

\noindent This paper extends the centered Magnus quadrature formulas of orders 2 to 8 in \cite{j28} to order 10 (residual error $O(t^{11})$ over the interval $[-t,t]$). The formula uses the centered Legendre moments $g_0,\ldots,g_4$ of the coefficient matrix and nine commutators. No formula with eight commutators was found. Of its 81 nonzero coefficients, 25 are small rationals and 56 are irrational algebraic numbers, defined as the isolated real solution of 70 polynomial equations with rational coefficients; they were refined by Newton's method to more than 190 digits, and the conditioning of the refinement is reported. The moments can be computed with five Gauss-Legendre nodes, with six Gauss-Lobatto nodes (five new samples per step), or with nine nodes that include the order-8 Newton-Cotes nodes. An embedded pair on seven Gauss-Kronrod nodes gives an order-6 estimate with three commutators and an order-10 estimate with seven more, whose difference serves as an error estimate.

\section*{Notation and order conditions}

The notation follows \cite{j28}. The Magnus exponent $\Omega(t)$ propagates the solution from $-t$ to $t$, and the centered Legendre moments are
\begin{equation}
g_{k}(t)=\int_{-t}^{t}P_{k}(\tau/t)\cdot A(\tau)\,d\tau\label{eq:gdef}
\end{equation}
(Eq. 35 of \cite{j28}). The moment $g_k$ is $O(t^{k+1})$ and has \emph{weight} $k+1$; $g_0$, $g_2$ and $g_4$ are odd functions of $t$, $g_1$ and $g_3$ are even, and $\Omega(t)$ contains only odd weights. The series of $\Omega$ in the moments, derived from the centered Magnus theorem (Eq. 34 of \cite{j28}) in exact rational arithmetic, is given through weight 7 in Eq. 41 of \cite{j28}. Its weight-9 part, which involves only $g_0,\ldots,g_4$, has 45 nonzero terms in the Lyndon basis:
\begin{equation}
\begin{aligned}
\Omega_{9}={}&-\tfrac{1}{2}[g_{3},g_{4}]+\tfrac{3}{28}[g_{0},[g_{2},g_{4}]]\\
&+\tfrac{7}{180}[[g_{0},g_{3}],g_{3}]-\tfrac{1}{28}[[g_{0},g_{4}],g_{2}]\\
&+\tfrac{9}{140}[g_{1},[g_{1},g_{4}]]-\tfrac{1}{4}[g_{1},[g_{2},g_{3}]]\\
&+\tfrac{1}{12}[[g_{1},g_{3}],g_{2}]-\tfrac{3}{280}[g_{0},[g_{0},[g_{1},g_{4}]]]\\
&-\tfrac{1}{72}[g_{0},[[g_{0},g_{3}],g_{2}]]+\tfrac{3}{280}[g_{0},[[g_{0},g_{4}],g_{1}]]\\
&+\tfrac{1}{280}[[g_{0},g_{1}],[g_{0},g_{4}]]-\tfrac{1}{30}[g_{0},[g_{1},[g_{1},g_{3}]]]\\
&-\tfrac{5}{168}[g_{0},[[g_{1},g_{2}],g_{2}]]-\tfrac{7}{120}[[g_{0},[g_{1},g_{3}]],g_{1}]\\
&-\tfrac{1}{72}[[g_{0},g_{2}],[g_{0},g_{3}]]+\tfrac{5}{168}[[g_{0},g_{2}],[g_{1},g_{2}]]\\
&-\tfrac{5}{168}[[[g_{0},g_{2}],g_{2}],g_{1}]+\tfrac{1}{40}[[[g_{0},g_{3}],g_{1}],g_{1}]\\
&-\tfrac{3}{56}[g_{1},[g_{1},[g_{1},g_{2}]]]+\tfrac{1}{1680}[g_{0},[g_{0},[g_{0},[g_{0},g_{4}]]]]\\
&-\tfrac{1}{720}[g_{0},[g_{0},[g_{0},[g_{1},g_{3}]]]]+\tfrac{1}{240}[g_{0},[g_{0},[[g_{0},g_{3}],g_{1}]]]\\
&+\tfrac{1}{240}[g_{0},[[g_{0},g_{1}],[g_{0},g_{3}]]]+\tfrac{1}{84}[g_{0},[g_{0},[g_{1},[g_{1},g_{2}]]]]\\
&+\tfrac{1}{84}[g_{0},[[g_{0},[g_{1},g_{2}]],g_{1}]]+\tfrac{5}{1008}[[g_{0},[g_{0},g_{2}]],[g_{0},g_{2}]]\\
&+\tfrac{1}{112}[g_{0},[[[g_{0},g_{2}],g_{1}],g_{1}]]-\tfrac{1}{180}[[g_{0},[g_{0},g_{3}]],[g_{0},g_{1}]]\\
&+\tfrac{1}{42}[[g_{0},g_{1}],[[g_{0},g_{2}],g_{1}]]+\tfrac{1}{56}[[[g_{0},g_{1}],g_{1}],[g_{0},g_{2}]]\\
&+\tfrac{3}{560}[[[[g_{0},g_{1}],g_{1}],g_{1}],g_{1}]+\tfrac{1}{4320}[g_{0},[g_{0},[g_{0},[g_{0},[g_{0},g_{3}]]]]]\\
&+\tfrac{1}{5040}[g_{0},[g_{0},[g_{0},[g_{0},[g_{1},g_{2}]]]]]+\tfrac{1}{1008}[g_{0},[g_{0},[g_{0},[[g_{0},g_{2}],g_{1}]]]]\\
&+\tfrac{1}{672}[g_{0},[g_{0},[[g_{0},g_{1}],[g_{0},g_{2}]]]]+\tfrac{1}{1120}[g_{0},[g_{0},[[[g_{0},g_{1}],g_{1}],g_{1}]]]\\
&-\tfrac{1}{672}[g_{0},[[g_{0},[g_{0},g_{2}]],[g_{0},g_{1}]]]-\tfrac{1}{504}[[g_{0},[g_{0},g_{1}]],[g_{0},[g_{0},g_{2}]]]\\
&-\tfrac{1}{1120}[g_{0},[[g_{0},g_{1}],[[g_{0},g_{1}],g_{1}]]]-\tfrac{1}{280}[[g_{0},[[g_{0},g_{1}],g_{1}]],[g_{0},g_{1}]]\\
&+\tfrac{1}{20160}[g_{0},[g_{0},[g_{0},[g_{0},[g_{0},[g_{0},g_{2}]]]]]]+\tfrac{1}{11200}[g_{0},[g_{0},[g_{0},[g_{0},[[g_{0},g_{1}],g_{1}]]]]]\\
&-\tfrac{1}{3360}[g_{0},[g_{0},[[g_{0},[g_{0},g_{1}]],[g_{0},g_{1}]]]]+\tfrac{1}{6720}[[g_{0},[g_{0},[g_{0},g_{1}]]],[g_{0},[g_{0},g_{1}]]]\\
&+\tfrac{1}{201600}[g_{0},[g_{0},[g_{0},[g_{0},[g_{0},[g_{0},[g_{0},g_{1}]]]]]]]
\end{aligned}\label{eq:series9}
\end{equation}

A formula of order 10 must reproduce the series at weights 3, 5, 7 and 9. The free Lie algebra generated by $g_0,\ldots,g_4$ has 1, 5, 16 and 48 basis elements at these weights (not counting the single letters $g_2$ and $g_4$, which occur in neither the series nor the formulas), so there are 70 \emph{order conditions} (Table~\ref{tab:conditions}, which extends Table 2 of \cite{j28}), compared with 0, 1, 6 and 20 at orders 2, 4, 6 and 8.

\begin{table}[H]
\centering\small
\caption{Order conditions and commutator count at order 10.}\label{tab:conditions}
\begin{tabular}{llllll}
\hline
Order & Residual & Moments & Conditions at weights 3, 5, 7, 9 & Total & Commutators \\ \hline
10 & $O(t^{11})$ & $g_0,\ldots,g_4$ & 1, 5, 16, 48 & 70 & 9 \\ \hline
\end{tabular}
\end{table}

These formulas were first derived for the conventional interval $[0,h]$ with $A$ expanded about the midpoint. The Legendre moments over $[0,h]$, with $P_k(2\tau/h-1)$, are the centered moments of Eq.~\ref{eq:gdef} with $h=2t$, and the two series agree term by term (as verified in \cite{j28} through weight 9), so the formulas carry over to the centered interval unchanged. All results below were re-verified against the centered series.

\section*{Number of commutators}

As in \cite{j28}, each formula is a sequence of nested commutators $C_i=[X_i,Y_i]$, where $X_i$ and $Y_i$ are linear combinations of the moments and earlier commutators and each has definite parity, followed by $\Omega=g_0+\sum_i c_i\cdot C_i$. The orders 4, 6 and 8 require 1, 3 and 6 commutators, the minimum numbers stated by Blanes, Casas and Ros \cite{bcr2002}; no order-10 formula has been published, as far as we know. The search at order 10 computed, for each pattern of operand parities, the maximum generic rank of the Jacobian of the 70 order conditions with respect to the free coefficients, and then searched for exact solutions with Levenberg-Marquardt iterations from random starting points (Table~\ref{tab:search}):

\begin{table}[H]
\centering\small
\caption{Search for the minimum number of commutators at order 10.}\label{tab:search}
\begin{tabular}{lllll}
\hline
Commutators & Patterns & Max. rank & Restarts & Result \\ \hline
8 & all 2,187 & 64 & 322 & residual $9\times10^{-4}$ \\
9 & all 6,561 & 69 & about 100 & exact, 3 families \\
10 & 3,000 sampled & 70 (26 patterns) & 11 & exact \\ \hline
\end{tabular}
\end{table}

\noindent A generic rank below 70 does not by itself exclude eight commutators, because the nine-commutator solutions lie on a special set where the Jacobian rank drops to 55-56; the failure at eight commutators is strong numerical evidence, not a proof.

\section*{The nine-commutator formula}

The operand parities of the chosen formula are, for $C_1$ to $C_9$, OE, OE, OO, OE, OO, OE, OO, OE and OE (O = odd, E = even); the commutators $C_3$, $C_5$ and $C_7$ are even and enter only as operands. With $C_i=[X_i,Y_i]$ and coefficients $x_1,\ldots,x_{81}$ (Appendix A), of which 25 are rational and 56 irrational,
\begin{equation}
\begin{aligned}
X_{1}&=x_{1}g_{0}+x_{2}g_{2}\\
Y_{1}&=x_{3}g_{1}+x_{4}g_{3}\\
X_{2}&=x_{5}g_{0}+x_{6}g_{2}+x_{7}C_{1}\\
Y_{2}&=x_{8}g_{1}+x_{9}g_{3}\\
X_{3}&=x_{10}g_{0}+x_{11}g_{2}\\
Y_{3}&=x_{12}g_{0}+x_{13}g_{2}+x_{14}g_{4}+x_{15}C_{1}\\
X_{4}&=x_{16}g_{0}+x_{17}g_{2}+x_{18}g_{4}+x_{19}C_{2}\\
Y_{4}&=x_{20}g_{1}+x_{21}g_{3}+x_{22}C_{3}\\
X_{5}&=x_{23}g_{0}+x_{24}g_{2}+x_{25}g_{4}\\
Y_{5}&=x_{26}g_{0}+x_{27}g_{2}+x_{28}g_{4}+x_{29}C_{4}\\
X_{6}&=x_{30}g_{0}+x_{31}g_{2}+x_{32}g_{4}+x_{33}C_{2}+x_{34}C_{4}\\
Y_{6}&=x_{35}g_{1}+x_{36}g_{3}+x_{37}C_{3}+x_{38}C_{5}\\
X_{7}&=x_{39}g_{0}+x_{40}g_{2}+x_{41}g_{4}+x_{42}C_{2}+x_{43}C_{4}+x_{44}C_{6}\\
Y_{7}&=x_{45}g_{0}+x_{46}g_{2}+x_{47}g_{4}+x_{48}C_{1}+x_{49}C_{4}+x_{50}C_{6}\\
X_{8}&=x_{51}g_{0}+x_{52}g_{2}+x_{53}g_{4}+x_{54}C_{1}+x_{55}C_{2}+x_{56}C_{4}+x_{57}C_{6}\\
Y_{8}&=x_{58}g_{1}+x_{59}g_{3}+x_{60}C_{3}+x_{61}C_{5}+x_{62}C_{7}\\
X_{9}&=x_{63}g_{0}+x_{64}g_{2}+x_{65}g_{4}+x_{66}C_{1}+x_{67}C_{2}+x_{68}C_{4}+x_{69}C_{6}+x_{70}C_{8}\\
Y_{9}&=x_{71}g_{1}+x_{72}g_{3}+x_{73}C_{3}+x_{74}C_{5}+x_{75}C_{7}\\
\Omega(t)&=g_{0}+x_{76}C_{1}+x_{77}C_{2}+x_{78}C_{4}+x_{79}C_{6}+x_{80}C_{8}+x_{81}C_{9}+O(t^{11})
\end{aligned}\label{eq:f10}
\end{equation}

The solution families have 16-20 free parameters beyond 24 rescaling freedoms. The member of Eq.~\ref{eq:f10} was chosen by minimizing the leading ($t^{11}$) error coefficients, with each moment weighted by its leading size, and then fixing as many coefficients as possible to simple rationals while keeping that weighted error within 25\% of the minimum; rescalings of the operands were then used to simplify the coefficients. Its weighted error is about 150 times smaller than for the member with the simplest coefficients, and the best points of the other families are within 6\% of it on the test problems. Attempts to make all coefficients rational failed: in every order of fixing that was tried, some coefficients became determined at irrational values.

\section*{The defining equations and root polishing}

With the 25 rational coefficients substituted, the order conditions become a system of polynomial equations in the vector $x$ of the 56 irrational coefficients. For each Lyndon word $\ell$ of odd weight 3 to 9 in the letters $g_0,\ldots,g_4$, the equation $E_\ell(x)=0$ states that the coefficient of the word $\ell$ in the expansion of the formula equals its coefficient in the series. Because a Lie polynomial is determined by its coefficients on the Lyndon words, the 70 equations are equivalent to the order conditions. They have rational coefficients, total degree at most 13 and 4951 terms in all; the weight-3 equation, for $\ell=g_0g_1$ and with $x_1=5/3$ substituted, is
\begin{equation}
\tfrac{1}{2} +\tfrac{5}{3}x_{3}x_{76} +x_{5}x_{8}x_{77} +x_{16}x_{20}x_{78} +x_{30}x_{35}x_{79} +x_{51}x_{58}x_{80} +x_{63}x_{71}x_{81}=0\label{eq:e1}
\end{equation}

\emph{Root polishing.} The coefficients were refined by Gauss-Newton iteration (Newton's method for the overdetermined, consistent system), $x\leftarrow x-(J^{T}J)^{-1}J^{T}E(x)$, with $E$ and its Jacobian $J$ evaluated from the exact rational coefficients in 320-digit decimal arithmetic. Convergence is quadratic. The largest $|E_\ell|$ at each iteration, starting from the coefficients rounded to 17 and to 30 significant digits, is shown in Table~\ref{tab:polish}.

\begin{table}[H]
\centering\footnotesize
\caption{Largest residual $|E_\ell|$ in Gauss-Newton polishing.}\label{tab:polish}
\begin{tabular}{lllllll}
\hline
Start & Iteration 0 & 1 & 2 & 3 & 4 & 5 \\ \hline
17 digits & $10^{-15}$ & $10^{-31}$ & $10^{-53}$ & $10^{-107}$ & $10^{-206}$ & $10^{-318}$ \\
30 digits & $10^{-28}$ & $10^{-57}$ & $10^{-105}$ & $10^{-210}$ & $10^{-318}$ &  \\ \hline
\end{tabular}
\end{table}

The final residual is at the level of the arithmetic precision, and the iterates agree with an earlier 150-digit solution to all of its digits. The values in Appendix A are rounded from this solution.

\emph{Conditioning.} The Jacobian $J=\partial E/\partial x$ ($70\times56$) has full column rank. Its singular values range from 11.5 to $5.9\times10^{-8}$, so its condition number is $\kappa=2.0\times10^{8}$; with each column scaled to unit norm, $\kappa=2.3\times10^{7}$. The two smallest singular values, $5.9\times10^{-8}$ and $3.6\times10^{-7}$, lie well below the rest (the next is $1.1\times10^{-5}$). Full rank means that the solution is isolated, so every $x_k$ is a real algebraic number. It also bounds the error of the polished coefficients by the final residual divided by the smallest singular value, about $10^{-310}$ (the solution was stored to 200 digits). In the opposite direction, rounding the coefficients to double precision changes the residuals by at most about $10^{-15}$, far below the $t^{11}$ truncation error at any useful step size, so the condition number does not limit the accuracy of the formula in double-precision use. Integer-relation tests (LLL on 140-digit values of $x_{2}$ and $x_{76}$) excluded minimal polynomials of degree 12 or less with integer coefficients below about $10^{10}$, so closed forms, if they exist, are large.

\section*{Moment quadrature}

By the criterion of \cite{j28} (moment-quadrature section), any rule of polynomial degree 9 gives moments accurate enough for order 10: $g_0$, which occurs alone, is then accurate to $O(t^{11})$, and $g_k$ for $k\ge1$ occurs only in terms of weight at least $2k+1$. Three such rules are given here. Each was verified by substituting its moments, expressed in the exact moments $g_0,\ldots,g_9$, into Eq.~\ref{eq:f10} and comparing every word through weight 10 with the series.

\emph{Gauss-Legendre, five nodes.} $g_k\approx 2t\cdot\sum_i w_i\cdot P_k(u_i)\cdot A(u_i\cdot t)$ (Eq. 42 of \cite{j28}, with $\sum_i w_i=1$), with the nodes and weights of Table~\ref{tab:gl5}, which extends Table 3 of \cite{j28}. The nodes are approximately $\pm0.5384693101$ and $\pm0.9061798459$.

\begin{table}[H]
\centering\small
\caption{Gauss-Legendre nodes and weights for order 10.}\label{tab:gl5}
\begin{tabular}{llll}
\hline
Nodes $n$ & Order & Nodes $u_i$ & Weights $w_i$ \\ \hline
5 & 10 & $0$ & $64/225$ \\
 &  & $\pm\frac{1}{3}\sqrt{5-2\sqrt{10/7}}$ & $(322+13\sqrt{70})/1800$ \\
 &  & $\pm\frac{1}{3}\sqrt{5+2\sqrt{10/7}}$ & $(322-13\sqrt{70})/1800$ \\ \hline
\end{tabular}
\end{table}

\emph{Gauss-Lobatto, six nodes.} The same form, with the nodes and weights of Table~\ref{tab:lo6}, which extends Table 5 of \cite{j28}. The interior nodes are approximately $\pm0.7650553239$ and $\pm0.2852315165$. The rule is exact for polynomials of degree 9. Because it includes both endpoints, the last sample of one step is the first of the next, so it needs five new samples of $A$ per step, like the Gauss-Legendre rule, and it also provides $A$ at the step boundaries. After a step is split, however, only the two endpoint samples can be reused, so the two half-steps need nine new samples. The weights must be applied to the products $P_k\cdot A$ as written; the alternative of integrating the interpolating polynomial through the six samples is exact only for $A$ of degree 5, which is not enough for $g_1,\ldots,g_4$.

\begin{table}[H]
\centering\small
\caption{Gauss-Lobatto nodes and weights for order 10.}\label{tab:lo6}
\begin{tabular}{llll}
\hline
Nodes $n$ & Order & Nodes $u_i$ & Weights $w_i$ \\ \hline
6 & 10 & $\pm1$ & $1/30$ \\
 &  & $\pm\sqrt{1/3+2\sqrt{7}/21}$ & $(14-\sqrt{7})/60$ \\
 &  & $\pm\sqrt{1/3-2\sqrt{7}/21}$ & $(14+\sqrt{7})/60$ \\ \hline
\end{tabular}
\end{table}

\emph{Nine nodes (modified Newton-Cotes).} With $u_i=-1,-\frac{5}{6},-\frac{2}{3},-\frac{1}{3},0,\frac{1}{3},\frac{2}{3},\frac{5}{6},1$, samples $a_i=A(u_i\cdot t)$, and weights obtained by integrating the interpolating polynomial against $P_k$,
\begin{equation}
\begin{aligned}
g_{0}&=\tfrac{t}{115500}\cdot(6235a_{0}+24576a_{1}+22440a_{2}+47025a_{3}+30448a_{4}\\
&\qquad+47025a_{5}+22440a_{6}+24576a_{7}+6235a_{8})\\
g_{1}&=\tfrac{t}{23100}\cdot(-1247a_{0}-4096a_{1}-2992a_{2}-3135a_{3}+3135a_{5}\\
&\qquad+2992a_{6}+4096a_{7}+1247a_{8})\\
g_{2}&=\tfrac{t}{1270500}\cdot(58295a_{0}+196608a_{1}-34320a_{2}-91575a_{3}-258016a_{4}\\
&\qquad-91575a_{5}-34320a_{6}+196608a_{7}+58295a_{8})\\
g_{3}&=\tfrac{t}{84700}\cdot(-3429a_{0}-8192a_{1}+9856a_{2}+11055a_{3}-11055a_{5}\\
&\qquad-9856a_{6}+8192a_{7}+3429a_{8})\\
g_{4}&=\tfrac{2t}{4129125}\cdot(69210a_{0}+80896a_{1}-282535a_{2}+1650a_{3}+261558a_{4}\\
&\qquad+1650a_{5}-282535a_{6}+80896a_{7}+69210a_{8})
\end{aligned}\label{eq:nine}
\end{equation}
The rule is exact for polynomials of degree 9 (degree 8 by interpolation, and degree 9 by symmetry). It needs eight new samples per step. Uniformly spaced nine-point Newton-Cotes nodes are not used, because of Runge's phenomenon: interpolation through equally spaced samples becomes unstable at high degree, and the nine-point Newton-Cotes weights alternate in sign, which amplifies rounding and data errors. (The $g_0$ weights above are all positive.) The nodes used here nevertheless lie on a uniform grid of spacing $1/6$ ($6u_i=0,\pm2,\pm4,\pm5,\pm6$), and when a step is split into two half-steps, the nodes of the two halves include all nine nodes of the original step, as with conventional Newton-Cotes rules. The split then needs only eight new samples instead of sixteen. The nodes also contain the seven equally spaced nodes of the order-8 Newton-Cotes rule of \cite{j28}, so an order-8 result is available from the same samples.

The leading quadrature error of $g_0$, $\int_{-t}^{t}A(\tau)\,d\tau-(\text{rule})=\kappa\cdot t^{11}\cdot A^{(10)}(0)$, has the coefficients in Table~\ref{tab:kappa}, which extends Table 4 of \cite{j28}; the Newton-Cotes column is for the nine-node rule. For the six-node Gauss-Lobatto rule $\kappa=-1/1031443875$, and the seven-node Gauss-Kronrod rule below is exact to degree 11, so its error enters only at $O(t^{13})$.

\begin{table}[H]
\centering\small
\caption{Leading quadrature-error coefficients $\kappa$ of $g_0$ at order 10.}\label{tab:kappa}
\begin{tabular}{lll}
\hline
Order & Gauss-Legendre $\kappa$ & Newton-Cotes $\kappa$ (nine nodes) \\ \hline
10 & $+1/1237732650$ & $-7/15588936000$ \\ \hline
\end{tabular}
\end{table}

\section*{Embedded pair on seven Gauss-Kronrod nodes}

This variant gives an order-6 estimate $M_6$ and an order-10 estimate $M_{10}$ from the same samples, and the difference $\|M_{10}-M_6\|$ estimates the error of $M_6$ and is a conservative bound for that of $M_{10}$. The seven Gauss-Kronrod nodes \cite{kronrod} contain the three Gauss-Legendre nodes (Table~\ref{tab:kronrod}).

\begin{table}[H]
\centering\small
\caption{Gauss-Kronrod and Gauss-Legendre nodes and weights.}\label{tab:kronrod}
\begin{tabular}{lll}
\hline
Node $u_i$ & Kronrod weight $w^K_i$ & Gauss weight $w^G_i$ \\ \hline
$0$ & $11008/48825$ & $4/9$ \\
$\pm\sqrt{3/5}$ & $6250/46557$ & $5/18$ \\
$\pm\frac{1}{3}\sqrt{5-6\sqrt{10/33}}\approx\pm0.4342437493$ & $\frac{2028807}{16036300}+\frac{1852389}{13745400}\sqrt{10/33}$ &  \\
$\pm\frac{1}{3}\sqrt{5+6\sqrt{10/33}}\approx\pm0.9604912687$ & $\frac{2028807}{16036300}-\frac{1852389}{13745400}\sqrt{10/33}$ &  \\ \hline
\end{tabular}
\end{table}

\noindent The Kronrod moments $g_k=2t\cdot\sum_{i=1}^{7}w^K_i\cdot P_k(u_i)\cdot A(u_i\cdot t)$ ($k=0,\ldots,4$) are exact for polynomials of degree 11 and behave like exact moments here. The Gauss moments $h_k=2t\cdot\sum w^G_i\cdot P_k(u_i)\cdot A(u_i\cdot t)$ ($k=0,1,2$, over the three Gauss nodes) are exact only to degree 5; in terms of exact moments,
\begin{equation}
\begin{aligned}
h_{0}&=g_{0}-\tfrac{429}{100}g_{6}+\tfrac{2431}{800}g_{8}+\ldots\\
h_{1}&=g_{1}-\tfrac{99}{50}g_{5}-\tfrac{22}{25}g_{7}+\tfrac{2223}{1000}g_{9}+\ldots\\
h_{2}&=g_{2}-\tfrac{27}{20}g_{4}-\tfrac{91}{1000}g_{6}-\tfrac{5151}{8000}g_{8}+\ldots
\end{aligned}\label{eq:hexp}
\end{equation}

\emph{Stage 1 (order 6).} From the Gauss samples alone, with three commutators,
\begin{equation}
\begin{aligned}
D_{1}&=\big[\tfrac{1}{2}h_{0}+h_{2},\ h_{1}\big]\\
D_{2}&=\big[h_{2}+\tfrac{1}{5}D_{1},\ h_{0}+D_{1}\big]\\
D_{3}&=\big[3h_{0}+11h_{2}+D_{1},\ h_{1}-\tfrac{5}{54}D_{2}\big]\\
M_{6}&=h_{0}-\tfrac{14}{5}D_{1}+\tfrac{3}{10}D_{3}+O(t^{7})
\end{aligned}\label{eq:stage1}
\end{equation}
which is an order-6 formula in the Gauss moments, equivalent to Eq. 47 of \cite{j28}.

\emph{Stage 2 (order 10).} Seven further commutators $C_i=[X_i,Y_i]$, $i=4,\ldots,10$, may use $g_0,\ldots,g_4$, $h_0,h_1,h_2$, $D_1,D_2,D_3$ and each other. Besides the 70 order conditions, they must cancel the terms containing $g_5,\ldots,g_8$ through weight 9 that enter through the $h_k$ (Eq.~\ref{eq:hexp}), which gives 78 conditions. With coefficients $y_1,\ldots,y_{96}$ (Appendix B), of which 35 are rational and 61 irrational,
\begin{equation}
\begin{aligned}
X_{4}&=y_{1}h_{0}+y_{2}h_{2}\\
Y_{4}&=y_{3}g_{1}+y_{4}g_{3}+y_{5}h_{1}+y_{6}D_{2}\\
X_{5}&=y_{7}g_{4}+y_{8}h_{0}+y_{9}h_{2}+y_{10}D_{1}+y_{11}D_{3}+y_{12}C_{4}\\
Y_{5}&=y_{13}g_{1}+y_{14}g_{3}+y_{15}h_{1}\\
X_{6}&=y_{16}g_{2}+y_{17}g_{4}+y_{18}h_{0}+y_{19}h_{2}+y_{20}D_{1}+y_{21}D_{3}+y_{22}C_{4}\\
&\qquad+y_{23}C_{5}\\
Y_{6}&=y_{24}g_{0}+y_{25}g_{2}+y_{26}g_{4}+y_{27}h_{0}+y_{28}h_{2}+y_{29}D_{1}+y_{30}D_{3}\\
&\qquad+y_{31}C_{4}\\
X_{7}&=y_{32}g_{0}+y_{33}g_{4}+y_{34}h_{0}+y_{35}h_{2}+y_{36}D_{1}+y_{37}C_{4}\\
Y_{7}&=y_{38}g_{1}+y_{39}g_{3}+y_{40}h_{1}+y_{41}D_{2}+y_{42}C_{6}
\end{aligned}\label{eq:stage2}
\end{equation}
\begin{equation}
\begin{aligned}
X_{8}&=y_{43}g_{0}+y_{44}g_{2}+y_{45}g_{4}+y_{46}h_{0}+y_{47}h_{2}+y_{48}C_{4}+y_{49}C_{7}\\
Y_{8}&=y_{50}g_{0}+y_{51}g_{2}+y_{52}g_{4}+y_{53}h_{0}+y_{54}h_{2}+y_{55}D_{3}+y_{56}C_{5}\\
&\qquad+y_{57}C_{7}\\
X_{9}&=y_{58}g_{2}+y_{59}g_{4}+y_{60}h_{0}+y_{61}h_{2}+y_{62}D_{1}+y_{63}D_{3}+y_{64}C_{4}\\
&\qquad+y_{65}C_{5}+y_{66}C_{7}\\
Y_{9}&=y_{67}g_{1}+y_{68}g_{3}+y_{69}h_{1}+y_{70}D_{2}+y_{71}C_{6}+y_{72}C_{8}\\
X_{10}&=y_{73}g_{0}+y_{74}g_{2}+y_{75}g_{4}+y_{76}h_{0}+y_{77}h_{2}+y_{78}D_{1}+y_{79}D_{3}\\
&\qquad+y_{80}C_{4}+y_{81}C_{5}+y_{82}C_{7}+y_{83}C_{9}\\
Y_{10}&=y_{84}g_{1}+y_{85}g_{3}+y_{86}h_{1}+y_{87}D_{2}+y_{88}C_{6}+y_{89}C_{8}\\
M_{10}&=g_{0}+y_{90}D_{1}+y_{91}D_{3}+y_{92}C_{4}+y_{93}C_{5}+y_{94}C_{7}+y_{95}C_{9}+y_{96}C_{10}+O(t^{11})
\end{aligned}\label{eq:stage2b}
\end{equation}
The large coefficients in $C_8$ multiply small combinations ($h_k-g_k$ and $g_4$); this is how stage 2 cancels the quadrature errors of the Gauss moments. With six new commutators (nine in all) the generic Jacobian rank reaches only 76, and 67 Levenberg-Marquardt restarts stalled at residual $1.1\times10^{-2}$; seven new commutators give exact solutions. The nodes must be the Kronrod nodes: a seven-node rule containing the Gauss nodes and the endpoints, $u=0,\pm\frac{1}{3},\pm\sqrt{3/5},\pm1$, is exact to degree 9, which suffices for Eq.~\ref{eq:f10}, but not for this pair, because stage 2 uses $g_4$ in weight-7 terms to cancel the $g_4$ component of $h_2$.

\emph{Root polishing.} The 61 irrational coefficients were polished as above, with the 78 conditions taken on the Lyndon words in $g_0,\ldots,g_8$. The Jacobian ($78\times61$) has full rank 61, so the solution is again isolated; its condition number is $3.2\times10^{10}$ ($1.9\times10^{8}$ with column scaling). Starting from the coefficients rounded to 20 significant digits, the largest residual fell from $10^{-18}$ to $10^{-36}$, $10^{-70}$, $10^{-139}$ and $10^{-208}$ in successive iterations, and the polished values differ from the 20-digit values by less than $4\times10^{-20}$ relative.

\section*{Cost per step}

Table~\ref{tab:cost}, which extends Table 6 of \cite{j28}, gives the cost per step of Eq.~\ref{eq:f10}; the Newton-Cotes column is for the nine-node rule. The embedded pair uses 10 commutators (20 matrix products) and 7 new samples per step.

\begin{table}[H]
\centering\small
\caption{Commutators, matrix products and samples of $A$ per step at order 10.}\label{tab:cost}
\begin{tabular}{llllll}
\hline
Order & Commutators & Matrix & Samples per step &  &  \\
 &  & products & Gauss- & Newton- & Gauss- \\
 &  &  & Legendre & Cotes (new) & Lobatto (new) \\ \hline
10 & 9 & 18 & 5 & 8 & 5 \\ \hline
\end{tabular}
\end{table}

\section*{Results}

The formulas were tested on the three test problems of \cite{j28} (random, commutator-heavy, and rapidly varying relative to its non-commutativity), with the same reference solution. Table~\ref{tab:results} lists the error in $\Omega(t)$ (Frobenius norm) for half-steps $t=0.1$ and $t=0.05$, and the error ratio between $t=0.1$ and $t=0.05$ for the random test. The table has the format of Table 7 of \cite{j28}. The order-10 rows give Eq.~\ref{eq:f10} with each moment rule (Newton-Cotes, 9 is the nine-node rule), followed by the embedded pair: $M_{10}$, $M_6$ and the error estimate $\|M_{10}-M_6\|$. The last row is the order-8 formula of \cite{j28} with four Gauss-Legendre nodes (Table 7 of \cite{j28}). The order-10 coefficients were read from the MATLAB file described below and used in double precision.

\begin{table}[H]
\centering\footnotesize\setlength{\tabcolsep}{4pt}
\caption{Errors in $\Omega(t)$ (Frobenius norm) for the three test problems.}\label{tab:results}
\begin{tabular}{lllllllll}
\hline
Order & Moments & Random &  & Comm.-heavy &  & Rapid &  & Ratio \\
 &  & $t=0.1$ & $t=0.05$ & $t=0.1$ & $t=0.05$ & $t=0.1$ & $t=0.05$ &  \\ \hline
10 & Gauss-Legendre, 5 & 1.1e-10 & 5.6e-14 & 9.2e-08 & 4.6e-11 & 2.3e-12 & 2.8e-15 & 2041 \\
 & Newton-Cotes, 9 & 1.1e-10 & 5.6e-14 & 9.2e-08 & 4.6e-11 & 1.1e-12 & 1.9e-15 & 2041 \\
 & Gauss-Lobatto, 6 & 1.1e-10 & 5.6e-14 & 9.2e-08 & 4.6e-11 & 2.5e-12 & 2.0e-15 & 2041 \\
 & Gauss-Kronrod, 7 & 1.1e-10 & 5.6e-14 & 9.2e-08 & 4.6e-11 & 3.2e-13 & 2.2e-15 & 2041 \\
 & exact moments & 1.1e-10 & 5.6e-14 & 9.2e-08 & 4.6e-11 & 3.2e-13 & 2.2e-15 & 2041 \\
10 (pair) & Gauss-Kronrod, 7 & 8.7e-11 & 4.3e-14 & 8.9e-08 & 4.5e-11 & 2.9e-13 & 2.1e-15 & 2033 \\
6 (pair) & Gauss-Legendre, 3 & 5.7e-07 & 4.5e-09 & 3.8e-05 & 3.1e-07 & 7.1e-08 & 5.5e-10 & 127 \\
estimate & $\|M_{10}-M_{6}\|$ & 5.7e-07 & 4.5e-09 & 3.8e-05 & 3.1e-07 & 7.1e-08 & 5.5e-10 & 127 \\
8 & Gauss-Legendre, 4 & 6.7e-09 & 1.3e-11 & 1.9e-06 & 3.8e-09 & 3.9e-10 & 7.7e-13 & 508 \\ \hline
\end{tabular}
\end{table}

Halving $t$ reduces the order-10 errors by factors approaching $2^{11}=2048$ and the $M_6$ errors by factors approaching $2^7=128$. (In the rapidly varying test at $t=0.05$ the order-10 errors reach the rounding level of about $2\times10^{-15}$.) In the first two tests all moment rules give the same errors to within 1\%. In the rapidly varying test the quadrature error of $g_0$ dominates: at $t=0.1$, exact or Kronrod moments are about 7 times more accurate than Gauss-Legendre moments, the nine-node rule is about 2 times more accurate, and the Lobatto rule is about 10\% less accurate. The embedded $M_{10}$ is as accurate as Eq.~\ref{eq:f10}, and $\|M_{10}-M_6\|$ tracks the error of $M_6$ to two or three digits. Compared with order 8, the order-10 formula costs three more commutators and one more sample per step, and at $t=0.1$ it is 20-170 times more accurate in these tests.

\section*{MATLAB coefficients}

The accompanying function magnus10coef.m, called as [x,y] = magnus10coef(), returns the coefficients $x_1,\ldots,x_{81}$ of Eq.~\ref{eq:f10} and $y_1,\ldots,y_{96}$ of Eqs.~\ref{eq:stage2} and \ref{eq:stage2b} as column vectors, in the order of Appendices A and B. Rational coefficients are written as exact fractions and irrational ones to 30 significant digits.

\section*{Code availability}

MATLAB implementations of the formulas (explde.m, magnus\_exp.m and magnus10coef.m), these documents and their LyX and LaTeX sources are available at \url{https://github.com/kjinnovation/explde} and on MATLAB File Exchange, \url{https://www.mathworks.com/matlabcentral/fileexchange/183695-explde}.

\section*{Acknowledgment}

This work was carried out with the assistance of Claude Opus 5.5 (Anthropic), a large language model. Under the author's direction, Claude derived the Magnus series in Legendre moments, searched for minimal-commutator formulas and optimized their leading error, and drafted portions of the text.

\section*{Appendix A: coefficients of the nine-commutator formula}

\noindent Tables~\ref{tab:xa}-\ref{tab:xc} list the coefficients $x_k$ of Eq.~\ref{eq:f10}; rationals are exact and irrational values are given to 30 significant digits. The position $X_i$, $g_j$ means the coefficient of $g_j$ in $X_i$.

\begin{table}[H]
\centering\small
\caption{Coefficients $x_1,\ldots,x_{27}$.}\label{tab:xa}
\begin{tabular}{lll}
\hline
 & Position & Value \\ \hline
$x_{1}$ & $X_{1}$, $g_{0}$ & $5/3$ \\
$x_{2}$ & $X_{1}$, $g_{2}$ & $0.0901217511054286179200006432152$ \\
$x_{3}$ & $Y_{1}$, $g_{1}$ & $0.132144147153367034682119532902$ \\
$x_{4}$ & $Y_{1}$, $g_{3}$ & $-1/5$ \\
$x_{5}$ & $X_{2}$, $g_{0}$ & $-0.472685390399902712406002183602$ \\
$x_{6}$ & $X_{2}$, $g_{2}$ & $3.23401783247220824419907849435$ \\
$x_{7}$ & $X_{2}$, $C_{1}$ & $1$ \\
$x_{8}$ & $Y_{2}$, $g_{1}$ & $0.161670971415242270251705225531$ \\
$x_{9}$ & $Y_{2}$, $g_{3}$ & $-0.525919202398768195923304595384$ \\
$x_{10}$ & $X_{3}$, $g_{0}$ & $-1.53079589966726923775289156382$ \\
$x_{11}$ & $X_{3}$, $g_{2}$ & $-0.910924396469683553350121373571$ \\
$x_{12}$ & $Y_{3}$, $g_{0}$ & $-19/20$ \\
$x_{13}$ & $Y_{3}$, $g_{2}$ & $-0.455192455511641898707317211627$ \\
$x_{14}$ & $Y_{3}$, $g_{4}$ & $-7/160$ \\
$x_{15}$ & $Y_{3}$, $C_{1}$ & $1/20$ \\
$x_{16}$ & $X_{4}$, $g_{0}$ & $0.999065729598492021595447046029$ \\
$x_{17}$ & $X_{4}$, $g_{2}$ & $-0.169685143102619456375901958438$ \\
$x_{18}$ & $X_{4}$, $g_{4}$ & $3.92051929889526396734420883485$ \\
$x_{19}$ & $X_{4}$, $C_{2}$ & $1$ \\
$x_{20}$ & $Y_{4}$, $g_{1}$ & $0.555796983434714384253066760376$ \\
$x_{21}$ & $Y_{4}$, $g_{3}$ & $-0.607334116992801989897737814155$ \\
$x_{22}$ & $Y_{4}$, $C_{3}$ & $1$ \\
$x_{23}$ & $X_{5}$, $g_{0}$ & $-8.07420917786156075183923102445$ \\
$x_{24}$ & $X_{5}$, $g_{2}$ & $-10.0495069461621075880489787299$ \\
$x_{25}$ & $X_{5}$, $g_{4}$ & $80/3$ \\
$x_{26}$ & $Y_{5}$, $g_{0}$ & $19/2$ \\
$x_{27}$ & $Y_{5}$, $g_{2}$ & $12.1017461973869000839204686308$ \\ \hline
\end{tabular}
\end{table}

\begin{table}[H]
\centering\small
\caption{Coefficients $x_{28},\ldots,x_{54}$.}\label{tab:xb}
\begin{tabular}{lll}
\hline
 & Position & Value \\ \hline
$x_{28}$ & $Y_{5}$, $g_{4}$ & $-31.4531547349911685094163487849$ \\
$x_{29}$ & $Y_{5}$, $C_{4}$ & $1/20$ \\
$x_{30}$ & $X_{6}$, $g_{0}$ & $0.275188457138167181875630064491$ \\
$x_{31}$ & $X_{6}$, $g_{2}$ & $-0.954326847042167934105954068326$ \\
$x_{32}$ & $X_{6}$, $g_{4}$ & $-0.784727265144758812578823778889$ \\
$x_{33}$ & $X_{6}$, $C_{2}$ & $1$ \\
$x_{34}$ & $X_{6}$, $C_{4}$ & $-1/4$ \\
$x_{35}$ & $Y_{6}$, $g_{1}$ & $1.61045764376873307258367256828$ \\
$x_{36}$ & $Y_{6}$, $g_{3}$ & $2.94807728420384234993029710640$ \\
$x_{37}$ & $Y_{6}$, $C_{3}$ & $-11$ \\
$x_{38}$ & $Y_{6}$, $C_{5}$ & $1$ \\
$x_{39}$ & $X_{7}$, $g_{0}$ & $-0.246517496492115672033127579776$ \\
$x_{40}$ & $X_{7}$, $g_{2}$ & $-2.90152105321322178449420125787$ \\
$x_{41}$ & $X_{7}$, $g_{4}$ & $-0.720231951037809406297947770664$ \\
$x_{42}$ & $X_{7}$, $C_{2}$ & $1$ \\
$x_{43}$ & $X_{7}$, $C_{4}$ & $-1/8$ \\
$x_{44}$ & $X_{7}$, $C_{6}$ & $-0.0834850346446151127595195646048$ \\
$x_{45}$ & $Y_{7}$, $g_{0}$ & $0.596648858778570464309441892471$ \\
$x_{46}$ & $Y_{7}$, $g_{2}$ & $-0.0848726298881828778243828476824$ \\
$x_{47}$ & $Y_{7}$, $g_{4}$ & $0.189439293642050895990550879245$ \\
$x_{48}$ & $Y_{7}$, $C_{1}$ & $-1$ \\
$x_{49}$ & $Y_{7}$, $C_{4}$ & $1/11$ \\
$x_{50}$ & $Y_{7}$, $C_{6}$ & $-2/9$ \\
$x_{51}$ & $X_{8}$, $g_{0}$ & $-1.22504203400172707389155574478$ \\
$x_{52}$ & $X_{8}$, $g_{2}$ & $2.93982839425824894217792481282$ \\
$x_{53}$ & $X_{8}$, $g_{4}$ & $1.89056056479215080737398254598$ \\
$x_{54}$ & $X_{8}$, $C_{1}$ & $1$ \\ \hline
\end{tabular}
\end{table}

\begin{table}[H]
\centering\small
\caption{Coefficients $x_{55},\ldots,x_{81}$.}\label{tab:xc}
\begin{tabular}{lll}
\hline
 & Position & Value \\ \hline
$x_{55}$ & $X_{8}$, $C_{2}$ & $-9/4$ \\
$x_{56}$ & $X_{8}$, $C_{4}$ & $0.493726863995418585916367639513$ \\
$x_{57}$ & $X_{8}$, $C_{6}$ & $0.143371545985261839634589032722$ \\
$x_{58}$ & $Y_{8}$, $g_{1}$ & $-1.10171588444014784576499503191$ \\
$x_{59}$ & $Y_{8}$, $g_{3}$ & $-1.33521791511463352749639989683$ \\
$x_{60}$ & $Y_{8}$, $C_{3}$ & $3$ \\
$x_{61}$ & $Y_{8}$, $C_{5}$ & $-0.300277475990164551753003485261$ \\
$x_{62}$ & $Y_{8}$, $C_{7}$ & $0.0422226321169889153042100461307$ \\
$x_{63}$ & $X_{9}$, $g_{0}$ & $-0.707315165639034053327374104654$ \\
$x_{64}$ & $X_{9}$, $g_{2}$ & $-1.00852021385304266703414054617$ \\
$x_{65}$ & $X_{9}$, $g_{4}$ & $0.522319462460385203547544946442$ \\
$x_{66}$ & $X_{9}$, $C_{1}$ & $-0.0395250285511817450646611919490$ \\
$x_{67}$ & $X_{9}$, $C_{2}$ & $1/2$ \\
$x_{68}$ & $X_{9}$, $C_{4}$ & $0.124884180986072604975597118288$ \\
$x_{69}$ & $X_{9}$, $C_{6}$ & $-0.501432555510506461635553953416$ \\
$x_{70}$ & $X_{9}$, $C_{8}$ & $0.335346238212012539516596859030$ \\
$x_{71}$ & $Y_{9}$, $g_{1}$ & $0.359077663794944462844741747862$ \\
$x_{72}$ & $Y_{9}$, $g_{3}$ & $0.837847882188203746637730745011$ \\
$x_{73}$ & $Y_{9}$, $C_{3}$ & $1/5$ \\
$x_{74}$ & $Y_{9}$, $C_{5}$ & $-0.0277515431965363360406864485263$ \\
$x_{75}$ & $Y_{9}$, $C_{7}$ & $-0.0751536499542879518128684277851$ \\
$x_{76}$ & $\Omega$, $C_{1}$ & $-0.594739381596017298987763177053$ \\
$x_{77}$ & $\Omega$, $C_{2}$ & $0.345400850217906386578044663611$ \\
$x_{78}$ & $\Omega$, $C_{4}$ & $-0.119041629440661798962689039945$ \\
$x_{79}$ & $\Omega$, $C_{6}$ & $0.195879050217731400879397915960$ \\
$x_{80}$ & $\Omega$, $C_{8}$ & $-0.219699011141684586146119923134$ \\
$x_{81}$ & $\Omega$, $C_{9}$ & $0.263055372971027419203395240270$ \\ \hline
\end{tabular}
\end{table}

\section*{Appendix B: coefficients of the embedded pair}

\noindent Tables~\ref{tab:ya}-\ref{tab:yc} list the coefficients $y_k$ of Eqs.~\ref{eq:stage2} and \ref{eq:stage2b}; rationals are exact and irrational values are given to 30 significant digits.

\begin{table}[H]
\centering\small
\caption{Coefficients $y_1,\ldots,y_{32}$.}\label{tab:ya}
\begin{tabular}{lll}
\hline
 & Position & Value \\ \hline
$y_{1}$ & $X_{4}$, $h_{0}$ & $0.410422029449033523794760374457$ \\
$y_{2}$ & $X_{4}$, $h_{2}$ & $-0.767184326790695144748574573849$ \\
$y_{3}$ & $Y_{4}$, $g_{1}$ & $1/2$ \\
$y_{4}$ & $Y_{4}$, $g_{3}$ & $-1.31232865221517550984653286316$ \\
$y_{5}$ & $Y_{4}$, $h_{1}$ & $-0.0639267136530839440218074034769$ \\
$y_{6}$ & $Y_{4}$, $D_{2}$ & $3/32$ \\
$y_{7}$ & $X_{5}$, $g_{4}$ & $0.911620344259285927256299010783$ \\
$y_{8}$ & $X_{5}$, $h_{0}$ & $-0.125735105194913262947328447707$ \\
$y_{9}$ & $X_{5}$, $h_{2}$ & $0.444046344967898289931020763694$ \\
$y_{10}$ & $X_{5}$, $D_{1}$ & $1/6$ \\
$y_{11}$ & $X_{5}$, $D_{3}$ & $-27/1760$ \\
$y_{12}$ & $X_{5}$, $C_{4}$ & $1/12$ \\
$y_{13}$ & $Y_{5}$, $g_{1}$ & $0.750097262567483160457495370942$ \\
$y_{14}$ & $Y_{5}$, $g_{3}$ & $-0.893372687955140272450281224406$ \\
$y_{15}$ & $Y_{5}$, $h_{1}$ & $0.0834337842172175186845577850221$ \\
$y_{16}$ & $X_{6}$, $g_{2}$ & $-1$ \\
$y_{17}$ & $X_{6}$, $g_{4}$ & $-1$ \\
$y_{18}$ & $X_{6}$, $h_{0}$ & $0.833384759757211049402703668837$ \\
$y_{19}$ & $X_{6}$, $h_{2}$ & $0.259828162285158389669148642865$ \\
$y_{20}$ & $X_{6}$, $D_{1}$ & $-1/3$ \\
$y_{21}$ & $X_{6}$, $D_{3}$ & $81/6160$ \\
$y_{22}$ & $X_{6}$, $C_{4}$ & $-0.170813649075048114297490034008$ \\
$y_{23}$ & $X_{6}$, $C_{5}$ & $1$ \\
$y_{24}$ & $Y_{6}$, $g_{0}$ & $-1/3$ \\
$y_{25}$ & $Y_{6}$, $g_{2}$ & $25/2$ \\
$y_{26}$ & $Y_{6}$, $g_{4}$ & $-18.6866839903490211274273667167$ \\
$y_{27}$ & $Y_{6}$, $h_{0}$ & $0.652737830562284394963761370734$ \\
$y_{28}$ & $Y_{6}$, $h_{2}$ & $-13.4931558741686603227232837885$ \\
$y_{29}$ & $Y_{6}$, $D_{1}$ & $-1/2$ \\
$y_{30}$ & $Y_{6}$, $D_{3}$ & $81/1760$ \\
$y_{31}$ & $Y_{6}$, $C_{4}$ & $-1/3$ \\
$y_{32}$ & $X_{7}$, $g_{0}$ & $-1/3$ \\ \hline
\end{tabular}
\end{table}

\begin{table}[H]
\centering\small
\caption{Coefficients $y_{33},\ldots,y_{64}$.}\label{tab:yb}
\begin{tabular}{lll}
\hline
 & Position & Value \\ \hline
$y_{33}$ & $X_{7}$, $g_{4}$ & $-0.595796530427459543546087044697$ \\
$y_{34}$ & $X_{7}$, $h_{0}$ & $0.544039246895247658107386805015$ \\
$y_{35}$ & $X_{7}$, $h_{2}$ & $0.0840492595773426301129101004253$ \\
$y_{36}$ & $X_{7}$, $D_{1}$ & $-1/8$ \\
$y_{37}$ & $X_{7}$, $C_{4}$ & $1/6$ \\
$y_{38}$ & $Y_{7}$, $g_{1}$ & $0.552499056229169796055542507726$ \\
$y_{39}$ & $Y_{7}$, $g_{3}$ & $0.0335436242400282608000330408492$ \\
$y_{40}$ & $Y_{7}$, $h_{1}$ & $-1.40554265878907087559830274908$ \\
$y_{41}$ & $Y_{7}$, $D_{2}$ & $-1/8$ \\
$y_{42}$ & $Y_{7}$, $C_{6}$ & $1/3$ \\
$y_{43}$ & $X_{8}$, $g_{0}$ & $-119/2$ \\
$y_{44}$ & $X_{8}$, $g_{2}$ & $-1387/2$ \\
$y_{45}$ & $X_{8}$, $g_{4}$ & $928.988353643688265914130123001$ \\
$y_{46}$ & $X_{8}$, $h_{0}$ & $60.0594208050479752625841092980$ \\
$y_{47}$ & $X_{8}$, $h_{2}$ & $693.692381734424602201762733442$ \\
$y_{48}$ & $X_{8}$, $C_{4}$ & $1$ \\
$y_{49}$ & $X_{8}$, $C_{7}$ & $3/2$ \\
$y_{50}$ & $Y_{8}$, $g_{0}$ & $-92$ \\
$y_{51}$ & $Y_{8}$, $g_{2}$ & $-914.900054948477505605246355089$ \\
$y_{52}$ & $Y_{8}$, $g_{4}$ & $1225.47233476213679589700346868$ \\
$y_{53}$ & $Y_{8}$, $h_{0}$ & $92.8363069307194257237209829503$ \\
$y_{54}$ & $Y_{8}$, $h_{2}$ & $915.140420340519614481104711094$ \\
$y_{55}$ & $Y_{8}$, $D_{3}$ & $243/1760$ \\
$y_{56}$ & $Y_{8}$, $C_{5}$ & $-5$ \\
$y_{57}$ & $Y_{8}$, $C_{7}$ & $6$ \\
$y_{58}$ & $X_{9}$, $g_{2}$ & $6$ \\
$y_{59}$ & $X_{9}$, $g_{4}$ & $-8.15969930921212323284616428682$ \\
$y_{60}$ & $X_{9}$, $h_{0}$ & $0.0748004497574308697823170967422$ \\
$y_{61}$ & $X_{9}$, $h_{2}$ & $-6.04813878778409449037469038704$ \\
$y_{62}$ & $X_{9}$, $D_{1}$ & $-1/20$ \\
$y_{63}$ & $X_{9}$, $D_{3}$ & $27/8800$ \\
$y_{64}$ & $X_{9}$, $C_{4}$ & $-0.0494619343106375461523035094178$ \\ \hline
\end{tabular}
\end{table}

\begin{table}[H]
\centering\small
\caption{Coefficients $y_{65},\ldots,y_{96}$.}\label{tab:yc}
\begin{tabular}{lll}
\hline
 & Position & Value \\ \hline
$y_{65}$ & $X_{9}$, $C_{5}$ & $0.100223590043548623475993317596$ \\
$y_{66}$ & $X_{9}$, $C_{7}$ & $-0.0800519646872135054608708988614$ \\
$y_{67}$ & $Y_{9}$, $g_{1}$ & $7.26294248076323644126729654658$ \\
$y_{68}$ & $Y_{9}$, $g_{3}$ & $0.789850680232171251637174378605$ \\
$y_{69}$ & $Y_{9}$, $h_{1}$ & $0.417013910658651610974633073895$ \\
$y_{70}$ & $Y_{9}$, $D_{2}$ & $-0.345137143611296325151139329226$ \\
$y_{71}$ & $Y_{9}$, $C_{6}$ & $10/3$ \\
$y_{72}$ & $Y_{9}$, $C_{8}$ & $-0.544501359611055470361583669568$ \\
$y_{73}$ & $X_{10}$, $g_{0}$ & $-1/5$ \\
$y_{74}$ & $X_{10}$, $g_{2}$ & $9.51291317660040269267562552719$ \\
$y_{75}$ & $X_{10}$, $g_{4}$ & $-12.8010538156900196830935421132$ \\
$y_{76}$ & $X_{10}$, $h_{0}$ & $0.158974968268125430175766915035$ \\
$y_{77}$ & $X_{10}$, $h_{2}$ & $-9.50378384438196258445870705920$ \\
$y_{78}$ & $X_{10}$, $D_{1}$ & $-1/10$ \\
$y_{79}$ & $X_{10}$, $D_{3}$ & $0.00709141050666411711074437892981$ \\
$y_{80}$ & $X_{10}$, $C_{4}$ & $0.00693770793595351213892558834982$ \\
$y_{81}$ & $X_{10}$, $C_{5}$ & $0.0619089624842630343080542131262$ \\
$y_{82}$ & $X_{10}$, $C_{7}$ & $-0.117839251134041588175713679959$ \\
$y_{83}$ & $X_{10}$, $C_{9}$ & $0.0595638129516088729514577667542$ \\
$y_{84}$ & $Y_{10}$, $g_{1}$ & $10$ \\
$y_{85}$ & $Y_{10}$, $g_{3}$ & $20.9859365144598103169261313693$ \\
$y_{86}$ & $Y_{10}$, $h_{1}$ & $-1.00602720808865272131737227031$ \\
$y_{87}$ & $Y_{10}$, $D_{2}$ & $0.466293441401308463352075347820$ \\
$y_{88}$ & $Y_{10}$, $C_{6}$ & $0.598435167563833017336934764848$ \\
$y_{89}$ & $Y_{10}$, $C_{8}$ & $0.102704991900605469294135573283$ \\
$y_{90}$ & $M_{10}$, $D_{1}$ & $-0.0745051862837507828645172568942$ \\
$y_{91}$ & $M_{10}$, $D_{3}$ & $0.00514753783338198345546041267870$ \\
$y_{92}$ & $M_{10}$, $C_{4}$ & $-0.358808845711455249074982220403$ \\
$y_{93}$ & $M_{10}$, $C_{5}$ & $0.520428513377888821232136671803$ \\
$y_{94}$ & $M_{10}$, $C_{7}$ & $-0.0751332166131541379289192802131$ \\
$y_{95}$ & $M_{10}$, $C_{9}$ & $-0.483085163978636981018871457461$ \\
$y_{96}$ & $M_{10}$, $C_{10}$ & $0.258603217397250949323128723863$ \\ \hline
\end{tabular}
\end{table}

\begin{thebibliography}{9}
\bibitem{j28} K. C. Johnson, ``The centered Magnus exponent and Magnus quadrature formulas of orders 2, 4, 6 and 8,'' Version 02-Nov-2026, magnus\_order\_2-8.pdf, \url{https://github.com/kjinnovation/explde}
\bibitem{bcr2002} S. Blanes, F. Casas, J. Ros, ``High order optimized geometric integrators for linear differential equations,'' \emph{BIT Numerical Mathematics} 42 (2002) 262-284.
\bibitem{kronrod} A. S. Kronrod, \emph{Nodes and Weights of Quadrature Formulas: Sixteen-Place Tables}, Consultants Bureau, New York, 1965.
\end{thebibliography}

\end{document}
